\chapter{Frostbite}

\section{Definition}
Frostbite is a cold-induced injury to tissues caused by freezing of water in cells and extracellular spaces. It most commonly affects exposed extremities (fingers, toes, ears, nose, cheeks). Severity ranges from superficial (first-degree) to deep tissue necrosis (fourth-degree).

\section{Pathophysiology}
Cold exposure causes vasoconstriction to preserve core temperature, reducing blood flow to extremities. At temperatures below 0 C (32 F), ice crystals form in extracellular fluid, causing osmotic shifts and cellular dehydration. Subsequent thawing and reperfusion triggers inflammatory cascade with microvascular thrombosis, edema, and tissue necrosis.

\section{Etiology}
\begin{itemize}
\item Prolonged exposure to freezing temperatures
\item Inadequate clothing or protective gear
\item High wind chill factor
\item Contact with cold surfaces or freezing metals
\item High altitude with decreased oxygen
\item Alcohol or drug intoxication (impaired judgment)
\item Peripheral vascular disease
\item Diabetes mellitus
\item Smoking (vasoconstriction)
\end{itemize}

\section{Indications for Evaluation}
\begin{itemize}
\item Cold, pale, and waxy skin
\item Numbness or loss of sensation in the affected area
\item Skin feels firm or frozen to the touch
\item After rewarming: pain, throbbing, burning sensation
\item Blisters or skin discoloration (red, blue, black)
\item Tissue necrosis or gangrene in severe cases
\end{itemize}

\section{Related Symptoms}
\begin{itemize}
\item Hypothermia (systemic cold exposure)
\item Chilblains (pernio, non-freezing cold injury)
\item Trench foot (immersion foot)
\item Raynaud phenomenon
\item Gangrene (tissue death)
\item Peripheral neuropathy
\end{itemize}

\section{Self-Care/Home Management}
\begin{itemize}
\item Move to a warm environment immediately
\item Remove wet or constrictive clothing
\item Immerse affected areas in warm water (37-39 C, 100-102 F) for 30-60 minutes
\item Do NOT use direct heat (heating pads, fire) - causes burns
\item Do NOT rub or massage frostbitten tissue
\item Do NOT break blisters
\item Elevate affected extremities
\end{itemize}

\section{Diagnostic Approach}
\begin{itemize}
\item Clinical assessment of depth and extent of injury
\item Rewarming in warm water and post-thaw assessment
\item Technetium-99m bone scan or angiography for deep injury
\item MRI or magnetic resonance angiography for tissue viability
\item Doppler ultrasound for blood flow assessment
\end{itemize}

\section{Treatment/Management Overview}
\begin{itemize}
\item Rapid rewarming in warm water (37-39 C) with analgesia
\item Tetanus prophylaxis if needed
\item Ibuprofen for inflammation
\item Topical aloe vera gel
\item Debridement of clearly necrotic tissue (delayed, 4-6 weeks)
\item Thrombolytic therapy with tPA for severe deep frostbite
\item Amputation for non-viable tissue (delayed until demarcation)
\end{itemize}
